\part*{Part III --- Placement}
\addcontentsline{toc}{part}{Part III --- Placement}

\bigskip
\noindent
\emph{Having a logic is not yet having an instrument. This part distinguishes
three things one can do with a formula --- synthesise a program from it, test a
program against it, or install it in a program as a guard --- and argues that the
third placement is what turns an observation discipline into something a system
can look through. It also introduces a fourth use, the formula as a revisable
hypothesis, which is what Part VII will need. Sources: \cite{ScientistNote} for
the assay and the hypothesis lattice, \cite{QuorumNote} for the reversal,
\cite{GsltOmnibus,GradedWhere} for the guard syntax.}
\bigskip

\section{Three placements}
\label{sec:placements}

Fix a generated logic $\Logic(\Sig,\Rw,\Coll,\V)$ and a formula $\varphi$ of it.
There are three essentially different things to do with $\varphi$, and they have
different costs, different prerequisites and different users.

\paragraph{Specification: $\vdash \varphi \triangleright P$.} From the formula,
produce a program realising it. This is proofs-as-programs, and for the logics of
Part II it is the direction with the most literature and the most open problems.
It requires a proof system, and for the coalition row there is not yet one: the
consensus-types development lists the extraction of a proof system as its central
gap, and four of its seven open problems are downstream of that gap
\cite{ConsensusTypes}.

\paragraph{Test: $P \models \varphi$.} Given a program, decide --- or, within a
budget, attempt to decide --- whether the formula holds. This requires only the
semantics, which the generator supplies with the connectives. This is what a model
checker does, and what an \emph{observer} does.

\paragraph{Guard: $\varphi$ inside the binder.} Install the formula in a
for-comprehension:
\begin{center}
\texttt{for(} $ptrn_{11} \leftarrow chan_{11}$ \texttt{\&} $\cdots$ \texttt{;}
$\cdots$ \texttt{;} $ptrn_{1n} \leftarrow chan_{1n}$ \texttt{ where } $\varphi$
\texttt{)}$P$
\end{center}
where \texttt{\&} joins within a group and \texttt{;} separates groups, and where
the preceding clauses bind the variables the condition may mention
\cite{GsltOmnibus}. Now the formula is evaluated when a communication is about to
occur, against the candidate data, and its value decides whether that
communication occurs.

The three are related but not interchangeable, and the differences are the
substance of this part.

\begin{table}[h]
\centering
\small
\renewcommand{\arraystretch}{1.3}
\begin{tabular}{@{}llll@{}}
\toprule
placement & judgement & needs & who uses it \\
\midrule
specification & $\vdash \varphi \triangleright P$ & a proof system & the designer \\
test          & $P \models \varphi$ & the semantics, plus a budget & the observer \\
guard         & $\varphi$ on the binder & per-match evaluation & the system itself \\
hypothesis    & $\varphi_{t} \rightsquigarrow \varphi_{t+1}$ & a lattice of formulae & the learner \\
\bottomrule
\end{tabular}
\caption{Four uses of one formula.}
\label{tab:placements}
\end{table}

\section{The arrow reverses, and the cheap direction is the useful one}
\label{sec:reversal}

The clearest instance of the difference is the coalition row, where the same logic
has been developed in both directions within this corpus.

The consensus-types development \cite{ConsensusTypes} reads the correspondence as
a term assignment. Its table maps participant to channel name, open set to channel
set, and the coalition modality at a participant to a \emph{join}; its right rule
introduces the join, its left rule introduces a send --- ``contribute local
evidence'' --- cut is communication, and cyclic proof is guarded recursion. The
computational content of the quorum-intersection lemma is that any two joins share
a channel carrying a well-typed message. It is a synthesis story: from a
specification of a consensus property, build a protocol.

A mortal observer wants the opposite. A hunting pride does not want to build a
herd; it wants to know which way the herd will break. Its question is
$P \models \varphi$, not $\vdash \varphi \triangleright P$.

\begin{proposition}[The observation reading sidesteps the proof theory]
\label{prop:sidestep}
Reading a generated logic as a hypothesis language about an already-running system
requires its semantics only. The proof-theoretic apparatus that synthesis requires
--- and which, for the coalition logic, does not yet exist --- is not on the
critical path.
\end{proposition}

This is the strongest strategic argument available for the whole
observation-discipline programme, and it is worth stating in general rather than
for one row. The generator hands you semantics for free, because the semantics
\emph{is} the lifting of constructors and contexts. It does not hand you a proof
system; proof systems have to be built, and building one is the expensive part of
adopting any new logic. An observer does not need one. So the range of logics that
are practically available for observation is far larger than the range that is
practically available for synthesis, and the dials of Part I can be turned freely
in the first setting long before the second catches up.

\begin{remark}[Where the guard sits]
The guard placement is on the observation side of this divide, not the synthesis
side. Evaluating $\varphi$ against a candidate message is a satisfaction question,
which is why a where-clause is cheap to add to a language and why it inherits
whatever fragment discipline the model checker was already under.
\end{remark}

\section{The guard, and what it is for}
\label{sec:guard}

\noindent\emph{In plain terms: a concurrent system has to decide, constantly,
which of several possible communications happens next. Ordinarily that decision is
made by the hardware. A guard moves it into the language, and makes it a question
about what the system can observe.}

A term of the form
\[
  x!(Q_{1}) \;\mid\; \texttt{for}(y \leftarrow x)P \;\mid\; x!(Q_{2})
\]
is a race. Two messages are available on one channel and one receipt is waiting;
exactly one of them will be consumed, and nothing in the calculus says which. In a
production implementation the answer is supplied by the physical race of actual
cores against a shared store. That is a resolution mechanism, but it is one the
language cannot speak about, cannot reason about, and cannot vary.

Installing a formula on the binder changes this. The guard is evaluated against
each candidate, and its value determines whether --- and, once graded, with what
propensity --- that candidate is taken. The non-determinism is still there, but it
is now resolved by \emph{something the system can say}.

\begin{principle}[The guard is where resolution becomes legible]
\label{prin:legible}
A logic in the specification placement describes a system. A logic in the guard
placement participates in it: the formula's value is the mechanism by which the
system's own non-determinism is resolved. Everything the observation discipline
can distinguish is thereby available to the scheduler, and nothing else is.
\end{principle}

Three consequences, each developed later.

\paragraph{Per-match evaluation.} The clause is evaluated once per candidate
match, not once per channel state. This is a design decision and it is the one
that makes the rest work: without it there is no map from candidates to values,
and hence nothing for Part V's weighing-versus-selecting distinction to act on.

\paragraph{Receipts are linear by default.} A linear receipt is consumed when it
fires; a persistent one is not. The default matters more than it looks. Under a
persistent receipt the term above is not a race at all: both messages are
eventually consumed, the receipt survives either way, and the two derivations
reconverge. It is \emph{linearity} that makes the surviving message a record of
which branch was taken, and therefore linearity on which the whole treatment of
racing in Parts IV and V rests \cite{GradedWhere}.

\paragraph{Perception must not be able to act.} If a formula is to read structure
without disturbing it, the reading must not itself be a communication that
consumes something the system needed. This is a real design constraint and it
bites in practice: in the board-game encoding, threat markers had to be given
their own name family, because a line signalling a threat on a square's own move
channel would rendezvous with that square's receipt and \emph{play the move}. The
same separation is forced in the herd model, where a scan implemented with an
ordinary receive rather than a peek would mean that the pride moves the zebra by
looking at it \cite{GameNote,QuorumNote}. Grades of access are separated by choice
of namespace, not by good intentions.

\section{The hypothesis placement}
\label{sec:hypothesis}

The fourth use of a formula is the one a learner makes: not to check a program
once, but to hold a guess and revise it.

\subsection{The assay}

The mortal-scientist development \cite{ScientistNote} makes an experiment out of
the guard placement directly. An \emph{assay} of a hypothesis $\varphi$ against a
specimen arriving on channel $p$ is a complementary pair of guarded receipts:
\begin{center}
\texttt{for(x <- p where }$\varphi(x)$\texttt{)\{ confirm \}}
\quad $\mid$ \quad
\texttt{for(x <- p where }$\lnot\varphi(x)$\texttt{)\{ refute \}}
\end{center}
Exactly one of three things then happens: the confirming branch fires, the
refuting branch fires, or neither does before the budget runs out. That
trichotomy --- confirm, refute, undetermined --- is the note's central
epistemological object, and the third outcome is not a modelling convenience: for
a mortal observer, \emph{refutation is budget-relative}. A step not yet enabled is
indistinguishable from a step never enabled, and telling them apart is exactly
what a finite purse cannot buy.

Two features of this construction will be picked up repeatedly.

First, it needs negation \emph{in the hypothesis logic}, and by
\S\ref{sec:adequacy} a generated logic supplies negation only if the target
specification does. The assay therefore carries an unadvertised dependency, which
Part IV discharges.

Second, the undetermined outcome has two sources that an observer cannot separate
from the outside. In the consensus setting, one is a property of the participant
--- stuck, faulty, Byzantine --- and the other is a property of the observer's
purse. A mortal scientist cannot distinguish a faulty participant from an
underfunded observation. Measured in the herd model: with a purse of $60$, $61.04\%$
of verdicts are undetermined, of which $56.68\%$ would have been decidable with a
richer purse; at a purse of $200$ or more the undetermined outcome disappears
entirely \cite{QuorumNote}.

\subsection{The lattice, and the shape of revision}

Hypotheses are formulae, and formulae are ordered by logical strength, so the
hypothesis space is a lattice --- the \emph{relaxation lattice}, obtained by
weakening a characteristic formula along its structure. It also carries a metric:
formulae agreeing to modal depth $n$ are within $2^{-n}$ of each other, which is
an ultrametric.

That geometry has consequences a designer should know before choosing a search
strategy.

\begin{itemize}[leftmargin=1.4em]
\item The balls are nested or disjoint, every point is a centre, and there is no
betweenness. \emph{There is no gradient.} Gradient methods are not ill-suited
here; they are ill-typed. What is available is tree search.
\item A revision move is a pair (locus, operation); its distance is the depth at
which one had to back up, and its direction is which sibling one moved to.
\item Cost is monotone in distance: a move of distance $2^{-n}$ invalidates every
confirmation obtained at depth $\ge n$. Backtracking is not free; it discards
evidence.
\item \textbf{Confinement.} A walk whose steps are all shorter than $2^{-n}$ never
leaves the ball of radius $2^{-n}$ it started in. A purely incremental learner
cannot escape its initial basin --- not because it is not clever enough, but for
metric reasons.
\end{itemize}

Confinement was measured, not merely proved. In the board-game study, the best
opening \emph{reverses} as theory depth increases: at one rung the centre opening
wins $.833$ against the corner's $.597$, while under the exactly-correct theory the
corner wins $.896$ against the centre's $.548$. Deepening the theory while holding
the opening fixed therefore makes things worse, $.833 \to .548$. The corner-opening
mutant invades at identical cost and is exactly the distance-$2^{-1}$ move that
Confinement denies an individual \cite{GameNote}.

\subsection{Why this part matters for the next one}

The picture just described has one degree of freedom and it is discrete. A learner
holds a formula; to change its behaviour it must change the formula; changing the
formula is a jump in a tree, priced by how far it had to back up, and small jumps
cannot leave the neighbourhood.

Part IV adds a second degree of freedom that is continuous and lives \emph{inside}
each ball. The consequences for learning are taken up in Part VII, but the
mechanism is entirely the subject of the next part: let the guard return a value
rather than a bit.
